\problemname{Game of the Throne}

\section{Problem}
In the ruthless world of Westeros, “The Game for The Iron Throne" is a dangerous game to play. 
One must navigate complex alliances and enemies. As the commander, you are tasked with strategizing. In order to claim victory, you and your alliance must strategically eliminate your enemies. 
However, your allies have specific preferences regarding which enemies they wish to kill, and as their leader, you must accommodate these preferences in your battle plans. For instance, we all know that Stannis harbors a not-so-secret wish to usurp his brothers’ claims, and Cersei? Well, she'd happily exterminate anyone who doesn't toast to her health with a glass of Dornish red. 
Each enemy belongs to a house, such as House Stark or House Targaryen. Since it's not always possible to eliminate all enemies from the same house at once, you must make calculated decisions about which members of each house to challenge strategically. 
Additionally, your alliance cannot always eliminate as many enemies as they would like. While Cersei may wish to kill everyone opposing her, there are practical limitations on how many targets you can engage at one time. 
Considering all these constraints, your task as a commander is to write a program that calculates the maximum number of enemies your alliance can kill. 
To successfully claim victory in the game for the Iron Throne, your strategy must ensure that you defeat more than 50\% of your enemies; failing to do so will require you to surrender.

\section{Input}
Vi mangler at finde ud af hvad vores MAX for a, e og h er.
The input begins with a line containing three integers \(a\), \(e\), and \(h\) separated by spaces. These represent the number of members in your alliance, the number of enemies, and the number of enemy houses, respectively \(1 \leq a, e \leq 100\), \(0 \leq h \leq e\). Members and enemies are uniquely identified by names.
The next \(a\) lines describe the preferences for each member of the alliance. Each line begins with the name of an alliance member followed by an integer \(k\), which indicates the number of enemies this member is willing to confront. The line continues with a list of names of enemies this member prefers to target.
Following this are \(h\) lines detailing the house affiliations and constraints. Each line starts with the name of an enemy house, followed by a list of enemy names affiliated with this house. The line ends with an integer \(r\), indicating the maximum number of members from this house that can be targeted. An enemy can belong to exactly one house.

\section{Output}
After calculating the maximum number of enemies your alliance can kill, the output is determined based on the percentage of enemies killed:
\begin{itemize}
    \item \textbf{If more than 50\% of the enemies have been eliminated}, the output will be "Dracarys" followed by the total number of enemies killed.
    \item \textbf{If 50\% or fewer of the enemies have been killed}, the output will be "When you play the game of thrones, you win or you die, therefore we surrender now" followed by the total number of enemies that could be killed.
\end{itemize}